\subsubsection*{Statement B.2q (Refuted): a constant cofinal Picard diagram over an imperfect field}
\begin{theorem}
\label{B:B2-cofinal-counterexample}
There is an imperfect field $k$ of characteristic $3$, a regular
function field $K/k$ of transcendence degree three, and two proper
geometrically normal models $G_1,G_2$ of $K/k$ with the following
property.  For every proper geometrically normal modification
$W\to G_i$, the geometric reduced identity component of
$\operatorname{Pic}_{W/k}$ is isomorphic to an elliptic curve
$E_i$.  The curves $E_1$ and $E_2$ are not geometrically isogenous.

Consequently there is no domination-cofinal collection of the
proper geometrically normal models of $K/k$ on which all the
$\operatorname{Pic}^{\tau}$ are isomorphic to one fixed group
scheme.  In particular a constant cofinal Picard diagram, with
comparison maps given by birational pullback, does not exist.
This does not refute the weak upper-cone formulation.
\end{theorem}

\begin{proof}
We give the construction and the Picard comparison separately.

\paragraph{The two elliptic factors and the generic section.}
Over $\mathbf F_3$ put
\[
 E_1:\ y^2=x^3-x,\qquad
 E_2:\ y^2=x^3+x^2+1,
 \qquad A=E_1\times E_2.
\]
Both plane cubics are smooth.  Write
$L=\mathcal O_{E_1}(1)\boxtimes\mathcal O_{E_2}(1)$.
This gives a projectively normal Segre embedding of $A$.
For completeness, each cubic section ring is its normal
homogeneous hypersurface coordinate ring: the affine cone is
regular off its vertex and satisfies $S_2$.  Their Segre product
is normal, since it is the invariant subring of their normal
tensor product for the action having weights $1$ and $-1$.
It is generated in degree one.  These arguments remain valid
after every field extension.

On each plane cubic, the six quadratic monomials
$X^2,XY,XZ,Y^2,YZ,Z^2$ form a basis of
$H^0(E_i,\mathcal O_{E_i}(2))$.  Let
$q_1,\ldots,q_{36}$ be their pairwise external products,
a basis of $H^0(A,L^2)$, and set
\[
 k=\mathbf F_3(a,c_1,\ldots,c_{36}),\qquad
 Q=\sum_{r=1}^{36}c_rq_r,
 \qquad l=k(a^{1/3}),
\]
where all the displayed variables are algebraically independent.
The divisor of $Q$ on $A_k$ is geometrically smooth.  One can
check the generic-section assertion directly: very ampleness of
$L^2$ makes vanishing of the value and the two first derivatives
at a geometric point three independent linear conditions on a
section.  The incidence of singular sections thus has dimension
at most $34$ in the product with the $35$-dimensional projective
linear system.  Its image is closed because $A$ is proper, and
does not contain the generic section.  The integer $36$ follows
from $h^0(E_i,\mathcal O_{E_i}(2))=6$ and the product formula.

\paragraph{The cone and its purely inseparable triple cover.}
Let $C$ be the projective cone over $(A_k,L)$, with cone coordinate
$W$ and vertex $o$.  Its homogeneous coordinate ring is
$S_A[W]$, where $S_A=\bigoplus_{n\ge0}H^0(A_k,L^n)$.
Define
\[
 M=\operatorname{Proj}
 \frac{S_A[W,V]}{(V^3-aW^3-WQ)}.
 \tag{B.2.7}
\]
Forgetting $V$ gives a finite flat map $M\to C$ of degree three,
with algebra
$\mathcal O_C\oplus\mathcal O_C(-1)\oplus\mathcal O_C(-2)$.
The unique point of $M$ over $o$ has residue field $l$; denote it
by $P$.

The threefold $M$ is geometrically integral and geometrically
normal.  Indeed, after adjoining $a^{1/3}$, the change
$V'=V-a^{1/3}W$ gives $V'^3=WQ$.  In the function field of the
cone the right side is not a cube, as its order in the variable
$W$ is one.  This proves geometric integrality.
The cone is geometrically normal and hence geometrically $S_2$;
the finite flat triple cover has zero-dimensional
Cohen--Macaulay fibres, so $M$ is geometrically $S_2$.
Away from $P$, the cone is smooth.  There the only singularities
of the displayed equation after an algebraic closure occur at
$W=V'=Q=0$.  Since $Q=0$ is smooth on the surface $A$, this is
a curve, of codimension two in $M$; the point $P$ has codimension
three.  Thus the geometric $R_1$ condition also holds, and
Serre's criterion proves geometric normality.  The commutative
algebra inputs for $S_2$ under field extension and flat maps are
\cite[Tags 045M, 0339 and 031S]{StacksB}.

\paragraph{The normal graph and its bad divisor.}
Let $p:R\to A_k$ be the ruled resolution of the cone, with
exceptional section $D_0$ of normal bundle $L^{-1}$, and let
$\rho:R\to C$ be the contraction.  Put
$H=\rho^*\mathcal O_C(1)$, so that $H=D_0+p^*L$.
Form the finite flat pullback
\[
                    \Gamma=M\times_C R.
\]
The incidence description of the ruled resolution embeds
$R$ as a closed subscheme of $C\times A_k$; hence $\Gamma$
is a closed subscheme of $M\times A_k$.
It is integral and isomorphic to $M$ away from $P$.
Its exceptional divisor is
\[
                       D=A_l.
\]
Locally along $D_0$, with parameter $t$ normal to the section
and a trivialization $q$ of $Q$, its equation is
\[
                         v^3=a+t^2q.                 \tag{B.2.8}
\]
The quotient by $t$ is the regular scheme $A_l$.
Since $t$ is a nonzerodivisor, this shows that $\Gamma$ is
regular at every point of $D$.  Away from $D$ it is the normal
scheme $M\setminus\{P\}$.  Thus $\Gamma$ is normal, and its
closed incidence embedding identifies it with the normal graph
of the rational map $M\dashrightarrow A_k$.

It is not geometrically normal.  At the generic point of
$D_{\bar k}$, equation (B.2.8) becomes
\[
                    v'^3=t^2q,\qquad q\ne0.
\]
This is a one-dimensional local domain of embedding dimension
two.  It is therefore not a discrete valuation ring and is not
normal.  A proper birational morphism between normal schemes is
an isomorphism over a neighbourhood of each codimension-one
point of the target.  Consequently no geometrically normal
proper model over $M$ can resolve $M\dashrightarrow A_k$:
such a resolution would factor through its closed graph
$\Gamma$ and would retain this geometrically nonnormal
codimension-one local ring.  The same obstruction persists
after every finite separable extension of $k$, because
$a^{1/3}$ remains outside such an extension.

\paragraph{Two geometrically normal partial graphs.}
For $i=1,2$, let $G_i$ be the normalization of the image of
\[
       \Gamma\longrightarrow M\times E_i,
\]
where the second component is the appropriate projection of
$A=E_1\times E_2$.  Normality of $\Gamma$ gives a proper
birational map $r_i:\Gamma\to G_i$.  Off $D$ it is an
isomorphism.  It contracts $D=A_l$ onto the curve $(E_i)_l$;
its fibres on $D$ are the other elliptic factor over $l$.
Indeed, the reduced image curve in $G_i$ is finite over
$(E_i)_l$, while the projection $A_l\to(E_i)_l$ has direct
image of its structure sheaf equal to $\mathcal O_{(E_i)_l}$.
Its Stein factorization therefore forces this finite
intermediate curve to be $(E_i)_l$ itself.
In particular $G_i\to M$ has no
exceptional divisor.

Every codimension-one point of $G_i$ therefore has a
neighbourhood isomorphic to an open of $\Gamma\setminus D$,
which is geometrically normal.  Thus $G_i$ is geometrically
$R_1$.  It is geometrically $S_2$ because it is normal:
field extensions have Cohen--Macaulay fibres, and $S_2$
ascends by the cited flat-map criterion.  Its function field
is the regular extension $k(M)/k$, so its geometric generic
fibre is integral.  Serre's criterion now proves that both
$G_i$ are geometrically normal.  They are projective, being
finite over closed subvarieties of $M\times E_i$.

\paragraph{The connected Picard scheme of the graph.}
For the degree-three map $\pi:\Gamma\to R$ we have
\[
 \pi_*\mathcal O_\Gamma
 =\mathcal O_R\oplus\mathcal O_R(-H)\oplus\mathcal O_R(-2H).
\]
The projective-bundle calculation gives
\[
 p_*\mathcal O_R=\mathcal O_A,\quad R^1p_*\mathcal O_R=0,
 \qquad Rp_*\mathcal O_R(-H)=0,
\]
and
\[
 p_*\mathcal O_R(-2H)=0,\qquad
 R^1p_*\mathcal O_R(-2H)=L^{-1}.
\]
The last identity also follows from relative duality, since
$\omega_{R/A}=\mathcal O_R(-2H)\otimes p^*L$.
As $H^0(A,L^{-1})=0$, we conclude that
\[
 H^1(\Gamma,\mathcal O_\Gamma)=H^1(A,\mathcal O_A),
 \qquad h^1(\Gamma,\mathcal O_\Gamma)=2.
\]
The map $\Gamma\to A$ has universal direct image of its
structure sheaf equal to $\mathcal O_A$, so pullback embeds
$A^\vee$ as a closed subgroup of $\operatorname{Pic}^0_{\Gamma/k}$.
This subgroup has dimension two, equal to the tangent dimension
of the ambient connected group scheme.  The latter is therefore
smooth of dimension two and equals $A^\vee$.

We also need a regular birational model with the same connected
Picard scheme.  The singularities of $\Gamma$ off $D$ are
\'etale locally the product of the surface
$xy=z^3$ with a smooth curve: replace $q$ locally by
$q+aW^2$ before using the equation.  Blowing up its reduced singular
locus resolves these singularities.  On the three charts of
the blowup of $(x,y,z)$ the strict transforms are respectively
$y_1=xz_1^3$, $x_1=yz_1^3$, and $z=x_1y_1$, hence smooth.
The exceptional surface-resolution fibre is a chain of two
projective lines of self-intersection $-2$.
For its reduced union $E$, the restriction of
$\mathcal O_E(-nE)$ has degree $n$ on each component for
$n\ge0$.  The normalization exact sequence at their one
intersection gives $H^1(E,\mathcal O_E(-nE))=0$.
Induction on infinitesimal neighbourhoods and formal functions
give $R^1$ of this surface resolution equal to zero; flat and
\'etale base change give the same conclusion for the local
products.  We thus obtain a proper regular model
$s:Y\to\Gamma$, unchanged near $D$, with
$R^1s_*\mathcal O_Y=0$.  Consequently
\[
                    \operatorname{Pic}^0_{Y/k}=A^\vee
                    \tag{B.2.9}
\]
by the same tangent-dimension argument.

Pullback along $r_i$ is a closed immersion of Picard schemes.
Its image in $\operatorname{Pic}^0_\Gamma=A^\vee$ contains
$E_i^\vee$, because $G_i$ maps to $E_i$.  On the other hand,
its restriction to every contracted fibre of
$A_l\to(E_i)_l$ is trivial.  On $A^\vee$ this condition is
exactly the vanishing of the other projection, after the
faithfully flat extension $l/k$.  It follows scheme-theoretically
that
\[
                    \operatorname{Pic}^0_{G_i/k}=E_i^\vee.
                    \tag{B.2.10}
\]

\paragraph{Every permitted domination stays on its elliptic branch.}
Let $W\to G_i$ be any proper geometrically normal modification.
Take the normalization $T$ of the main graph between $W$ and
the fixed regular model $Y$.  The map $T\to Y$ has normal
source and regular target, so B.1g identifies
$\operatorname{Pic}^{\tau}_T$ with
$\operatorname{Pic}^{\tau}_Y$.  Pullback from $W$ is a closed
immersion.  We obtain compatible closed immersions
\[
       E_i^\vee\ \subseteq\ \operatorname{Pic}^0_{W/k}
              \ \subseteq\ A^\vee.                 \tag{B.2.11}
\]
Suppose the middle group has dimension two.  A closed subgroup
of the abelian surface with that dimension and connected
underlying space is the whole abelian surface.  Hence the
second immersion is an isomorphism over $k$.

To turn this equality into the needed morphism, we may make a
finite separable field extension and choose a rational point
in the common smooth birational locus of $T,Y,W$.  Their
Albanese torsors then acquire compatible origins.  For proper
geometrically integral schemes the Albanese variety is dual to
the maximal abelian subvariety of the Picard scheme
\cite[Theorem 3.2 and Remark 3.3, pp.~12--13 of the author
manuscript]{C_ACMV2026}.  The two pullback isomorphisms to
$\operatorname{Pic}^0_T=A^\vee$ therefore dualize to
isomorphisms of Albanese varieties.  The morphism
$T\to Y\to A$ factors through a morphism $W\to A$, up to
the translation fixed by the chosen origin.  This extends the
original rational map $M\dashrightarrow A$ on $W$.
It contradicts the bad-divisor obstruction, which survives
this separable extension.  Thus the middle group in (B.2.11)
has dimension one.  After an algebraic closure its reduced
identity component is exactly $E_i^\vee$.

\paragraph{The two branches cannot have one constant value.}
We finish by checking the distinction between the elliptic
curves without an unverified Hasse-invariant calculation.
The chord-and-tangent law
\cite[Chapter I, Section 3, Theorem 3.1, pp.~25--27]{MilneEC2021B}
says that a finite point of a smooth plane cubic is killed by
$3$ precisely when its tangent is a flex.  For
$y^2=x^3+a_2x^2+a_4x+a_6$ in characteristic three and $y\ne0$,
substitution of the tangent line shows that this condition is
$m^2=a_2$, where $m$ is its slope.  On $E_1$ we have
$a_2=0$ and $m=-1/(2y)\ne0$, so there is no such finite point;
the points with $y=0$ are of order two.  On $E_2$ we have
$m=x/y$, so the flex condition gives $y^2=x^2$, hence
$x^3+1=0$.  The two solutions are $(-1,1)$ and $(-1,-1)$.
Thus
\[
 \#E_1[3](\bar k)=1,\qquad \#E_2[3](\bar k)=3.
\]
Surjectivity of multiplication by $3$ on geometric points
gives respectively $1$ and $3^r$ points killed by $3^r$.
If the two curves were isogenous, dualizing if necessary would
give an isogeny $E_2\to E_1$.  It would kill every geometric
$3^r$-torsion point of $E_2$, placing $3^r$ distinct points in
one fixed finite kernel for every $r$, which is impossible.
The two curves are not geometrically isogenous.
The accompanying exact program verifies the cubic smoothness,
tangent substitution and geometric flex equations; it keeps
the nonreduced cube equation distinct from its two reduced
finite flex points.

Finally, a domination-cofinal collection must contain a model
dominating $G_1$ and a model dominating $G_2$.  An isomorphism
between their $\operatorname{Pic}^{\tau}$ would induce an
isomorphism of their geometric reduced identity components,
contradicting the preceding calculation.  This proves the
asserted failure of the strong cofinal-diagram formulation.
\end{proof}

\begin{remark}[Scope]
We have not asserted that the entire torsion Picard group is
constant on either upper cone.  Finite infinitesimal subgroup
schemes in the other elliptic factor may still occur in
(B.2.11).  The reduced identity components already distinguish
the two cofinal requirements and suffice for the counterexample.
\end{remark}
